% ============================================================
%  MÓDULO 08 — O Ecossistema: configuração, agentes, skills,
%  hooks, plugins, especificações e diagnóstico. Parte VII.
% ============================================================
\chapter{Como o Ecossistema Trabalha: Agentes, Instruções e Verificações}

\roteirodidatico
{Para realizar uma tarefa, o operador pode combinar uma ficha de agente, instruções reutilizáveis e verificações automáticas. O capítulo explica a diferença entre essas peças.}
{\emph{Agente} descreve uma configuração de trabalho; \emph{skill} reúne instruções para um procedimento; \emph{plugin} adiciona integração; \emph{hook} reage a um evento; \emph{teste} verifica uma condição.}
{Siga uma tarefa desde a configuração até a seleção de instruções, execução de ferramentas e verificação. Consulte os arquivos de origem e observe o que o runtime realmente carrega.}
{A contagem de registros descreve o catálogo; não mede ativação, execução bem-sucedida ou qualidade. Examine também versão, permissões, chamadas observáveis e resultados.}
{O que o ambiente carregou nesta execução, qual componente foi acionado e que evidência mostra o efeito do hook ou da verificação?}

\casoguiado{perfil, instrução e verificação}
{Um perfil descreve um papel; uma skill pode orientar como revisar um manuscrito; um hook pode reagir a uma ação; um teste pode conferir uma regra específica. Esses itens não são sinônimos e nenhum deles, isoladamente, prova a qualidade da revisão. Para reconstruir a execução, registre o perfil carregado, a instrução usada, a ação observada e o critério aplicado à saída.}

\begin{descricao}
\textbf{Descrição — o que faltava no livro.} Antes, o capítulo apresentava
nomes como \emph{skill}, \emph{hook}, \emph{runtime}, \emph{MCP} e
\emph{spec} sem explicar com clareza o que cada um significa nem como afeta o
trabalho do operador. Também não separava com precisão um recurso documentado
de um recurso configurado, disponível ou efetivamente testado. Esta revisão
define esses termos, percorre o caminho da ficha até a execução, apresenta
exemplos calculados e aponta limites encontrados no código.
\end{descricao}

\section{Configuração: das fichas ao aplicativo}

\elemento{Do catálogo à configuração em uso}{\metainfo{Rota}{opencode.json, .opencode/} \hfill \metainfo{Fonte}{\nodep{integrations/opencode_cli.py}}}

\begin{descricao}
O \pth{opencode.json} é um arquivo de configuração: nele o aplicativo OpenCode
encontra instruções, agentes, skills, conexões MCP e comandos que o projeto
oferece. MCP (Model Context Protocol) é o protocolo que permite a um aplicativo
solicitar ferramentas a um servidor separado. O arquivo é produzido por
\pth{integrations/opencode_cli.py}, a rotina que lê fontes do projeto e escreve
uma configuração para o aplicativo. Por isso, o catálogo e os documentos são as
fontes de edição; o JSON é o resultado gerado.
\end{descricao}

\begin{conceito}
Três palavras aparecem ao longo do capítulo. \textbf{CLI} é a interface de
linha de comando: o operador digita um comando em um terminal. \textbf{JSON} é
um formato de texto estruturado usado para guardar opções e listas de
configuração. \textbf{Runtime} significa o programa em execução, com os recursos
que conseguiu carregar naquela sessão. Assim, ``estar no JSON'' e ``estar
disponível no runtime'' são situações diferentes.
\end{conceito}

\begin{resultado}
Em palavras simples: \pth{instructions} aponta para regras de trabalho;
\pth{agent} lista perfis disponíveis; \pth{skills} aponta para instruções
reutilizáveis; \pth{mcp} descreve conexões com servidores de ferramentas;
\pth{command} define comandos nomeados; e \pth{permission} define o que cada
rota pode fazer. O diagnóstico local de 29/09/2026 encontrou 216 entradas de
agente na configuração e 7 conexões MCP declaradas. Uma entrada declarada não
garante que o agente ou servidor esteja instalado ou pronto para uso.
\end{resultado}

\begin{flowch}[Regeneração da configuração]{documentos $\to$ opencode.json}
\matrix[matrix of nodes,name=cfg,row sep=5mm,column sep=4mm,nodes={fn},ampersand replacement=\&]{
 |[fns]|AGENTS.md + catálogo + docs \\[2mm]
 |[fns]|\nodep{integrations/opencode_cli.py} \\[2mm]
 |[fns]|\nodep{compila (modos, permissões, MCP)} \\[2mm]
 |[fns]|escreve opencode.json \\[2mm]
 |[fns]|CLI carrega ecossistema \\ \\
};
\draw[seta] (cfg-1-1)--(cfg-2-1);
\draw[seta] (cfg-2-1)--(cfg-3-1);
\draw[seta] (cfg-3-1)--(cfg-4-1);
\draw[seta] (cfg-4-1)--(cfg-5-1);
\draw[seta,plumAccent,dashed] (cfg-5-1.west) to[out=180,in=180,looseness=1.4]
  node[left=2pt,fill=papel,inner sep=1pt,text=plumAccent,font=\scriptsize]{repete}
  (cfg-1-1.west);
\end{flowch}

\begin{descricao}
\textbf{Como ler o diagrama.} Uma ficha do catálogo é a descrição original; o
compilador transforma essas descrições em JSON; o aplicativo lê o JSON durante
sua inicialização. Chamamos o JSON de \emph{derivado} porque ele foi produzido a
partir de outras fontes. Se alguém o editar diretamente, a próxima geração pode
apagar a edição. A regra prática é alterar a ficha ou o documento de origem e
depois regenerar o arquivo.

\textbf{Por que essa regra existe.} Manter uma fonte principal evita que duas
versões discordantes disputem qual delas vale. A ideia é parecida com guardar
uma planilha de dados e gerar um gráfico a partir dela: se o gráfico estiver
errado, corrige-se a planilha ou o procedimento que o produz. Os princípios
FAIR são orientações para que dados sejam encontráveis, acessíveis,
interoperáveis e reutilizáveis; também destacam a importância de identificar
as fontes de informação \citref{WILKINSON, M. D. et al. The FAIR Guiding Principles for scientific data management and stewardship. \emph{Scientific Data}, v.~3, 160018, 2016}{10.1038/sdata.2016.18}{O princípio de encontrabilidade ajuda o leitor a localizar a fonte que deve ser reutilizada.}{Findability is the first step in the reuse of data}{A localização é o primeiro passo para reutilizar dados.}{WILKINSON et al., 2016, p.~160018}.
\end{descricao}

\begin{metodo}
Para atualizar um agente: (1) edite sua ficha Markdown no catálogo; Markdown é
um formato de texto simples com títulos e listas; (2) execute
\pth{python3 -m integrations.opencode_cli} para gerar novamente a configuração;
(3) rode \pth{python3 -m marceloclaro.cli doctor} para verificar se as partes
principais do ambiente respondem. O diagnóstico identifica condições locais;
ele não testa automaticamente toda resposta produzida por um agente.
\end{metodo}

\clearpage
\section{Os agentes: da ficha à tarefa}

\elemento{Do catálogo à execução de uma tarefa}{\metainfo{Fonte}{agents/catalog/} \hfill \metainfo{Configuração local}{216 entradas}}

\begin{descricao}
Um \emph{agente} é um perfil de trabalho para o modelo: combina instruções,
identidade e capacidades para tratar certo tipo de pedido. A ficha original
fica em \pth{agents/catalog/}. No início do arquivo há um bloco chamado
\emph{frontmatter}: campos de metadados escritos entre delimitadores, como nome,
descrição, versão, categoria e exemplos. Metadados são dados que descrevem
outros dados; neste caso, explicam ao sistema para que serve o perfil.

O compilador transforma fichas em entradas reconhecidas pelo OpenCode. O
\emph{orquestrador} é o componente que coordena tarefas e escolhe caminhos de
execução. Quando usa o Blackboard, publica uma tarefa no quadro compartilhado;
agentes elegíveis podem se oferecer para realizá-la. Blackboard significa aqui
um quadro de tarefas consultável pelos componentes, não uma lista de espera de
pessoas. O diagnóstico de 29/09/2026 encontrou 216 entradas configuradas, mas
esse número não prova que todas estejam instaladas, disponíveis ou capazes de
concluir qualquer tarefa.
\end{descricao}

\begin{resultado}
As fichas cobrem, entre outros, pesquisa, escrita científica, saúde, análise de
código, produção editorial, apresentações e ferramentas externas. Por exemplo,
uma ficha de revisão de código descreve que tipo de arquivo examina e que saída
deve produzir. Essa descrição ajuda a localizar um perfil possível; a qualidade
do trabalho ainda precisa ser avaliada pela tarefa e por seus critérios.
\end{resultado}

\begin{flowch}[Ciclo de vida do agente]{catálogo $\to$ registro $\to$ operação}
\matrix[matrix of nodes,name=ag,row sep=5mm,column sep=4mm,nodes={fn},ampersand replacement=\&]{
 |[fns]|ficha no catálogo \\[2mm]
 |[fns]|compilador gera perfil para delegação \\[2mm]
 |[fns]|configuração (modo e instruções) \\[2mm]
 |[fns]|registro no quadro de tarefas \\[2mm]
 |[fns]|voluntariado e execução \\[2mm]
 |[fns]|reflexão e evolução \\ \\
};
\draw[seta] (ag-1-1)--(ag-2-1);
\draw[seta] (ag-2-1)--(ag-3-1);
\draw[seta] (ag-3-1)--(ag-4-1);
\draw[seta] (ag-4-1)--(ag-5-1);
\draw[seta] (ag-5-1)--(ag-6-1);
\draw[seta,plumAccent,dashed] (ag-6-1.west) to[out=180,in=180,looseness=1.4]
  node[left=2pt,fill=papel,inner sep=1pt,text=plumAccent,font=\scriptsize]{atualiza}
  (ag-1-1.west);
\end{flowch}

\begin{descricao}
\textbf{O que acontece em cada etapa.} A ficha declara o propósito do agente;
o compilador cria a entrada de configuração; o orquestrador ou o aplicativo
seleciona uma rota; se houver uma tarefa publicada no Blackboard, um agente
elegível pode aceitá-la; por fim, o resultado volta ao componente que solicitou
o trabalho. Há duas situações que não devem ser confundidas: \emph{não
selecionado} significa que outro candidato foi escolhido ou que faltou uma
condição; \emph{reprovado} significa que o agente executou, mas o resultado não
atendeu ao critério.

\textbf{Por que mostrar falhas no diagrama.} Um fluxo que mostra apenas sucesso
esconde o que o operador deve fazer quando o trabalho não é aceito ou falha. Na
literatura, agentes são descritos por sua capacidade de reagir, agir em busca de
objetivos e interagir com outros componentes. Aqui, reação significa responder
a uma tarefa ou evento; objetivo significa perseguir o resultado pedido dentro
das regras; interação significa devolver estado e resultado ao orquestrador.
Esses conceitos ajudam a explicar o modelo, mas não certificam a qualidade de
uma implementação específica \citref{WOOLDRIDGE, M.; JENNINGS, N. R. Intelligent agents: theory and practice. \emph{The Knowledge Engineering Review}, v.~10, n.~2, p.~115-152, 1995}{10.1017/S0269888900008122}{Os autores discutem agentes como tema importante da inteligência artificial e da ciência da computação.}{The concept of an agent has become important in both artificial intelligence and mainstream computer science}{O conceito de agente tornou-se importante tanto na inteligência artificial quanto na ciência da computação.}{WOOLDRIDGE; JENNINGS, 1995, p.~115--152}.
\end{descricao}

\begin{niveis}
\textbf{Intuição:} a ficha de um especialista explica que trabalho ele pode
receber; a ficha, sozinha, não prova que ele esteja de plantão.\\
\textbf{Implementação:} o compilador lê os metadados e gera a configuração.\\
\textbf{Termos:} A2A significa comunicação entre agentes. Registro A2A é o
cadastro de capacidades para que o roteador possa localizar candidatos; não é
uma avaliação independente de competência.
\end{niveis}

\clearpage
\section{Skills: instruções reutilizáveis}

\elemento{Instruções reutilizáveis para tarefas específicas}{\metainfo{Fonte}{.opencode/skills/} \hfill \metainfo{Política}{varia por skill}}

\begin{resultado}
Uma \emph{skill} é um conjunto reutilizável de instruções, exemplos e, às vezes,
programas auxiliares. Ela ensina um procedimento para um tipo de tarefa; não é
por si só um modelo novo nem a garantia de que a tarefa será executada. Exemplos
de funções presentes neste repositório:\par
\smallskip
{\footnotesize
\begin{tabularx}{\textwidth}{lY}
\toprule
Skill & Papel \\
\midrule
\pth{customize-opencode} & orienta a edição da configuração do aplicativo. \\
\pth{deploy-estatico-github-pages} & descreve as etapas para publicar um site estático. \\
\pth{gemini-cli} & explica como chamar o programa de linha de comando do Gemini sem interface gráfica. \\
\pth{gemini-notebooklm-bridge} & orienta a comunicação entre o Gemini e o NotebookLM. \\
\pth{goose-cli} & descreve como encaminhar tarefas ao programa Goose. \\
\pth{haystack-rag} & orienta RAG (do inglês, geração aumentada por recuperação): busca trechos em documentos e acrescenta contexto a uma resposta. \\
\pth{medico-virtual-supremo} & orienta apoio clínico sujeito a revisão humana. \\
\pth{mirofish-offline} & descreve uma simulação social executada localmente. \\
\pth{nano-orchestration} & divide a produção de manuscritos longos em etapas. \\
\pth{pesquisa-artigo-qualis-a1} & orienta pesquisa e redação com critérios auditáveis; não promete classificação. \\
\pth{pesquisador-universal-marcelo-claro} & descreve etapas de pesquisa reproduzíveis. \\
\pth{plandex-cli} & orienta o planejamento de mudanças de código com Plandex. \\
\pth{prestacao-contas-nota-dez} & orienta a conferência de diligências escolares. \\
\pth{reasonix-cli} & descreve a chamada do programa Reasonix. \\
\pth{smoke-test-dom-js} & orienta um teste rápido de JavaScript que usa elementos de página. \\
\bottomrule
\end{tabularx}\par}
\end{resultado}

\begin{arquitetura}
Cada skill costuma ter um arquivo \pth{SKILL.md}, o manual que explica quando e
como seguir aquele procedimento. Algumas incluem também scripts ou código de
apoio. A configuração pode exigir invocação explícita: isso significa que a
skill só é usada quando o operador a pede ou quando o fluxo a seleciona de forma
autorizada. Uma marca como \pth{disable-model-invocation: true} impede que o
modelo a escolha por conta própria; o orquestrador ainda pode ler suas
instruções conforme a política definida. Essa regra evita iniciar rotinas
sensíveis sem que a rota prevista tenha sido escolhida.
\end{arquitetura}

\clearpage
\section{Plugins}

\elemento{Código que amplia o aplicativo OpenCode}{\metainfo{Fonte}{.opencode/plugins/}}

\begin{resultado}
Um \emph{plugin} é um complemento de código carregado pelo aplicativo para
acrescentar comportamento. Diferentemente de uma skill, que explica um
procedimento, o plugin pode executar funções em resposta a chamadas do
aplicativo. Os plugins encontrados neste projeto são:\par
\smallskip
{\footnotesize
\begin{tabularx}{\textwidth}{lY}
\toprule
Plugin & Papel \\
\midrule
\pth{deploy-guards.ts} & verifica condições de publicação, como tamanho e disponibilidade das páginas. \\
\pth{litert-lm-provider.ts} & adapta o serviço local LiteRT-LM para uma interface de provedor compatível. \\
\bottomrule
\end{tabularx}\par}
\end{resultado}

\begin{limite}
\textbf{Provider}, ou provedor, é o componente que conecta o aplicativo a um
serviço de modelo. \textbf{Fallback} é uma alternativa usada quando o caminho
preferido não está disponível. Neste projeto, a conexão que segue o formato de
chamadas da interface OpenAI é descrita primeiro no \pth{opencode.json}; o plugin LiteRT-LM oferece um
caminho complementar. Portanto, plugin, provider e skill ocupam papéis
diferentes: código de extensão, conexão com modelo e instruções de tarefa.
\end{limite}

\clearpage
\section{Hooks: as guardas de segurança}

\elemento{Verificações automáticas ligadas a eventos}{\metainfo{Fonte}{.opencode/hooks/}}

\begin{resultado}
Um \emph{hook} é um programa executado quando acontece um evento definido, como
antes de uma publicação. O evento aciona a verificação sem depender de alguém
lembrar de seguir uma recomendação escrita. Guardas encontradas no diretório
\pth{.opencode/hooks/}:\par
\smallskip
{\footnotesize
\begin{tabularx}{\textwidth}{lY}
\toprule
Hook & Função \\
\midrule
\pth{budget\_guard.sh} & interrompe a rota quando o limite configurado de uso é excedido. \\
\pth{credential\_guard.sh} & procura credenciais, como chaves de acesso, antes que sejam expostas. \\
\pth{js\_smoke\_dom.sh} & faz uma verificação rápida de JavaScript que manipula elementos de uma página. \\
\bottomrule
\end{tabularx}\par}
\end{resultado}

\begin{metodo}
O teste rápido de página funciona assim:
\begin{enumerate}[leftmargin=1.4em,itemsep=1pt]
  \item identifica alterações em JavaScript que procuram ou mudam elementos da página;
  \item cria um \emph{stub}, isto é, um substituto pequeno que imita as funções necessárias do navegador;
  \item executa o trecho de JavaScript com Node.js, um programa que executa JavaScript fora do navegador;
  \item verifica se o seletor encontrou o elemento esperado. Se receber um valor vazio e o código tentar usá-lo, o erro é exibido;
  \item registra falha para que a rotina de publicação não trate o teste como aprovado.
\end{enumerate}
\end{metodo}

\begin{niveis}
\textbf{Intuição:} um hook é como uma cancela automática: ela confere a condição
no momento em que a passagem é solicitada.\\
\textbf{Diferença para uma revisão manual:} o teste executado observa um
comportamento simples; uma leitura do código apenas procura padrões escritos.
Ambas são úteis, mas respondem a perguntas diferentes.\\
\textbf{Limite:} passar nesse teste rápido não prova que o site inteiro funciona
em todos os navegadores ou dispositivos.
\end{niveis}

\clearpage
\section{O registro de especificações e a saúde do ecossistema}

\elemento{Especificações, testes e diagnóstico do ambiente}{\metainfo{Fonte}{specs/, marceloclaro/cli.py}}

\begin{resultado}
Uma \emph{spec} (especificação) descreve o comportamento esperado de uma parte do
sistema e os critérios que permitem avaliá-lo. O arquivo da spec registra a
intenção; o teste executável verifica alguns desses critérios no código. Um
\emph{ciclo de evolução} é outro tipo de registro: anota uma mudança e seu
contexto. Uma spec não é uma execução, e um ciclo não é prova de que todos os
testes passaram.

Na verificação local de 29/09/2026, o comando \pth{doctor} carregou 136
especificações formais e leu 457 ciclos de evolução. O registro passou a conter
458 ciclos após a anotação editorial R632 e 459 após esta revisão R633. O total de arquivos em
\pth{specs/} pode incluir documentos históricos e materiais auxiliares; por
isso, não deve ser confundido com o número de specs formais carregadas.
\end{resultado}

\begin{funcao}
O comando \pth{python3 -m marceloclaro.cli doctor} é um diagnóstico: verifica
condições locais como leitura das specs, integridade do registro de evolução e
disponibilidade de algumas integrações. No retrato de 29/09/2026, 18 das 20
verificações passaram, duas emitiram avisos e nenhuma falhou. Um aviso indica
algo ausente ou opcional que merece atenção; não equivale automaticamente a uma
falha geral. Por exemplo, uma ferramenta externa pode não estar instalada.

\emph{Memória metacognitiva} significa registros que o sistema guarda sobre seu
próprio trabalho anterior, como reflexões ou avaliações. O \pth{helpdesk} é uma
ajuda guiada; o \pth{menu} reúne opções da linha de comando. Esses comandos têm
funções distintas: diagnóstico, orientação e navegação.
\end{funcao}

\begin{flowch}[Diagnóstico de saúde]{verificações locais antes do trabalho}
\matrix[matrix of nodes,name=dg,row sep=5mm,column sep=4mm,nodes={fn},ampersand replacement=\&]{
 |[fns]|doctor (diagnóstico local) \\[2mm]
 |[fns]|especificações e histórico de mudanças \\[2mm]
 |[fns]|memória de reflexões disponível \\[2mm]
 |[fns]|servidores de ferramentas e modelos acessíveis \\[2mm]
 |[fns]|ambiente saudável? \\[2mm]
 |[fns]|prosseguir trabalho \\ \\
};
\draw[seta] (dg-1-1)--(dg-2-1);
\draw[seta] (dg-2-1)--(dg-3-1);
\draw[seta] (dg-3-1)--(dg-4-1);
\draw[seta] (dg-4-1)--(dg-5-1);
\draw[seta] (dg-5-1)--(dg-6-1);
\draw[seta,plumAccent,dashed] (dg-5-1.west) to[out=180,in=180,looseness=1.4]
  node[left=2pt,fill=papel,inner sep=1pt,text=plumAccent,font=\scriptsize]{dor}
  (dg-1-1.west);
\end{flowch}

\begin{descricao}
\textbf{Como interpretar o resultado.} O diagnóstico percorre verificações e
classifica cada uma como aprovada, falha ou aviso. Uma aprovação quer dizer que
aquela verificação local passou naquele momento. Um aviso relata uma condição
que não bloqueou o diagnóstico, como uma ferramenta opcional ausente. Uma falha
indica que uma condição necessária não foi satisfeita. Essas categorias ajudam
o operador a escolher o próximo passo; elas não avaliam, por si sós, se uma
resposta de pesquisa ou um texto está correto.

\textbf{Por que registrar incertezas.} Em psicologia, metacognição inclui
monitorar e regular o próprio pensamento. Neste sistema, usamos a expressão de
modo operacional: guardar sinais sobre verificações, execuções e limitações
anteriores. Isso permite distinguir ``não está instalado'' de ``foi executado e
falhou'', uma distinção necessária para agir sobre o diagnóstico \citref{FLAVELL, J. H. Metacognition and cognitive monitoring: a new area of cognitive-developmental inquiry. \emph{American Psychologist}, v.~34, n.~10, p.~906-911, 1979}{10.1037/0003-066X.34.10.906}{Flavell descreve o monitoramento ativo e a regulação dos processos cognitivos como parte da metacognição.}{Metacognition refers, among other things, to the active monitoring and consequent regulation and orchestration of these processes}{A metacognição inclui o monitoramento ativo e a regulação dos processos cognitivos.}{FLAVELL, 1979, p.~906--911}.
\end{descricao}

\section{Resumo e exercícios}

\begin{resultado}
\textbf{O que você deve ter aprendido:} uma ficha descreve um agente; o
compilador gera entradas de configuração; uma skill ensina um procedimento; um
plugin amplia o aplicativo com código; um hook executa uma verificação quando
ocorre um evento; uma spec registra comportamento esperado; e um teste verifica
um critério observável. O diagnóstico mostra condições locais, não certifica
toda a capacidade do ecossistema.

\textbf{Exercício.} Imagine que você alterou uma ficha de agente. Explique com
suas palavras: onde está a fonte original, qual comando atualiza a configuração,
o que o diagnóstico consegue verificar e o que ainda exigiria executar uma
tarefa real.
\end{resultado}

% ============================================================

%  Gerado por gen_mod14_ativacao.py (R613) — seção de ativação e cálculo
%  para o módulo mod-08-ecossistema.tex. Edições manuais são preservadas nas próximas
%  execuções (marcador impede reescrita).
% ============================================================

\section[Leitura guiada]{Leitura guiada — O que o inventário mostra}

\elemento{Leitura guiada — Declaração não é execução}{\metainfo{ID}{N-08} \hfill \metainfo{Estilo}{inventário e evidência}}

\begin{descricao}
\textbf{A pergunta.} Se a configuração lista 216 agentes, isso quer dizer que
216 agentes estão prontos para trabalhar? Não necessariamente. O número conta
entradas declaradas no arquivo de configuração. Para saber se um agente pode
ser usado, é preciso verificar se sua ficha foi carregada, se as ferramentas
necessárias existem e se a rota consegue executá-lo.

\textbf{Uma descrição também importa.} A descrição é o resumo que informa ao
roteador que tipos de tarefa combinam com um agente. Uma frase como ``ajuda com
pesquisa'' pouco distingue um especialista em revisão bibliográfica de um
especialista em estatística. Uma descrição mais útil especifica a entrada que o
agente espera, o trabalho que realiza e o tipo de resultado que entrega. Isso
melhora a informação disponível para escolher; não garante que a resposta será
correta.
\end{descricao}

\begin{definicao}
\textbf{Quatro estados que não devem ser confundidos.}
\begin{itemize}[leftmargin=1.3em,itemsep=1pt]
  \item \textbf{Documentado:} existe uma ficha de texto com nome e descrição.
  \item \textbf{Configurado:} a ficha aparece no arquivo que o aplicativo lê.
  \item \textbf{Disponível:} o processo atual consegue carregar o perfil e suas
        dependências.
  \item \textbf{Executado e avaliado:} uma tarefa foi realmente processada e
        seu resultado foi comparado com critérios definidos.
\end{itemize}
\end{definicao}

\begin{metodo}
\textbf{Exemplo de cálculo.} Suponha, apenas para ilustrar a conta, que um
catálogo tenha 20 fichas e que 18 delas apareçam na configuração. A cobertura
de registro seria $18/20=0{,}90$, ou 90\%. Isso mede a correspondência entre
fichas e configuração; não mede quantos agentes estão ativos nem se produzem
bons resultados.
\begin{enumerate}[leftmargin=1.4em,itemsep=2pt]
  \item O numerador é o número de fichas que têm uma entrada correspondente.
  \item O denominador é o total de fichas que se pretende comparar.
  \item O quociente responde apenas à pergunta ``quantas estão registradas?''.
  \item Para medir execução, seria necessário definir tarefas de avaliação,
        executá-las e registrar quantos resultados satisfizeram seus critérios.
\end{enumerate}
\end{metodo}

\begin{resultado}
\textbf{O que sabemos no retrato local.} Em 29/09/2026, o diagnóstico contou 216
entradas de agentes na configuração. Esse retrato não fornece, por si só, o
número de agentes que passaram por uma tarefa de avaliação. Não se deve chamar
essa contagem de ``216 agentes testados'' ou ``216 agentes prontos'': são
afirmações diferentes e exigem evidências diferentes.
\end{resultado}

\begin{resultado}
\textbf{Como verificar.} Compare fichas do catálogo com entradas geradas na
configuração e anote quais não têm correspondência. Depois escolha uma tarefa
representativa, execute-a pela rota documentada e guarde entrada, saída e
critério de avaliação. O primeiro passo confere registro; o segundo observa
comportamento. Um não substitui o outro.
\end{resultado}

\begin{aviso}
\textbf{Limite de validade.} A contagem não mede capacidade. Um agente pode estar
descrito e configurado, mas não ser selecionado, não ter uma ferramenta
necessária ou produzir uma resposta inadequada. Para falar em desempenho,
informe qual tarefa foi executada, sob quais condições e com quais critérios.
\end{aviso}

\begin{resultado}
\textbf{Exercícios.} (1) Em suas palavras, qual a diferença entre um agente
documentado, configurado e testado? (2) Se 18 de 20 fichas tiverem registro
correspondente, calcule a cobertura. O que essa porcentagem ainda não permite
concluir? (3) Dê um exemplo de uma descrição que ajude a escolher entre dois
agentes parecidos.
\end{resultado}

% ATIVACAO-R613
\section[Economia interna]{Economia interna: caução, penalidade e preço}
\elemento{Como são usados os termos econômicos}{\metainfo{Ciclo}{R613} \hfill \metainfo{Módulo}{mod-08-ecossistema.tex}}

\begin{processo}
\textbf{Token Economy} significa, neste projeto, um mecanismo interno de
contabilização de unidades associado a tarefas. A palavra \emph{token} aqui não
significa criptomoeda nem dinheiro real. O orquestrador pode consultar regras
desse mecanismo durante determinadas rotas; a presença do código não quer dizer
que toda chamada do sistema passe por ele.
\end{processo}

\begin{resultado}
\begin{description}[leftmargin=2.6cm,style=nextline]
  \item[Stake (caução)] unidades reservadas quando um agente aceita uma tarefa.
        A reserva representa um compromisso interno; não é depósito financeiro.
  \item[Slashing (desconto)] redução da reserva ou do saldo quando uma regra
        definida registra uma falha. A regra precisa indicar o evento e o cálculo.
  \item[Fee market (preço da tarefa)] regra que atribui custos internos segundo
        fatores como prioridade. É uma forma de ordenar ou limitar recursos, não
        um preço cobrado em moeda.
  \item[Trust score (indicador de confiança)] medida interna usada em algumas
        regras de encaminhamento. Um valor abaixo de 0,3 pode exigir supervisão
        em rotas configuradas; esse número não é uma probabilidade estatística
        de que o agente esteja certo.
\end{description}
\end{resultado}

\begin{descricao}
\textbf{Por que usar limites.} Uma fila pode priorizar pedidos urgentes, mas
dispõe de tempo e capacidade finitos. Do mesmo modo, um agente escolhe entre as
opções que conhece e pode executar dentro do prazo e das permissões recebidas.
Em economia e ciência da decisão, esse quadro é chamado de \emph{racionalidade
limitada}: não se supõe que o decisor conheça todas as alternativas nem possa
calcular tudo. O histórico pode informar regras de encaminhamento, mas não
transforma um indicador interno em prova de qualidade \citref{SIMON, H. A. A behavioral model of rational choice. \emph{The Quarterly Journal of Economics}, v.~69, n.~1, p.~99-118, 1955}{10.2307/1884852}{Simon propõe estudar escolhas compatíveis com a informação e a capacidade computacional realmente disponíveis ao decisor.}{Broadly stated, the task is to replace the global rationality of economic man with a kind of rational behavior that is compatible with the access to information and the computational capacities that are actually possessed by organisms}{Em termos amplos, a proposta é estudar um comportamento compatível com o acesso à informação e as capacidades computacionais realmente disponíveis.}{SIMON, 1955, p.~99}.
\end{descricao}

\begin{resultado}
\textbf{Cálculo com os valores definidos no módulo.} O código declara saldo
inicial de 100 unidades internas, caução mínima de 1 unidade, recompensa de
20\% em caso de sucesso e desconto de 50\% da caução em caso de falha. Se uma
caução de 2 unidades for registrada, o sucesso devolve $2+0{,}2\times2=2{,}4$
unidades; a falha devolve $2\times(1-0{,}5)=1$ unidade. A conta explica a regra
implementada na classe econômica; não transforma essas unidades em dinheiro.

Para a taxa, a prioridade normal custa 1 unidade quando não há congestionamento;
prioridade alta aplica multiplicador 2, portanto custa 2 unidades nessa mesma
condição. Dez tarefas abertas acrescentam 10\%: a taxa normal passa a
$1\times1\times1{,}1=1{,}1$ unidade. O número de tarefas abertas e o nível de
prioridade entram na fórmula descrita em \pth{economy/token\_economy.py}.
\end{resultado}

\begin{aviso}
\textbf{Limite.} Essas unidades e regras organizam decisões internas do software.
Elas não comprovam que a entrega seja verdadeira nem substituem revisão humana,
testes adequados ou validação externa. Para dizer que um agente recebeu desconto,
é necessário consultar a rota executada e o registro da tarefa correspondente.
\end{aviso}

\clearpage
\begin{aviso}
\textbf{Divergência encontrada na integração.} A classe
\pth{TokenEconomy.commit} recebe o parâmetro \pth{stake}, mas a chamada atual do
orquestrador usa \pth{amount=2.0}. Na versão inspecionada, Python rejeita essa
chamada porque o nome não corresponde à assinatura; o bloco captura a exceção e
segue adiante. A taxa pode já ter sido debitada pela instrução anterior, mas a
caução não é registrada nem há estorno nesse bloco. Antes de afirmar que um
desconto ocorreu, confira a assinatura e o registro da tarefa. A divergência está em
\pth{marceloclaro/orchestrator.py} e \pth{economy/token\_economy.py}.
\end{aviso}

\noindent\textit{Exemplo.} Se uma rota admite apenas cinco candidatos por limite de tempo e permissões, ela pode comparar os cinco que passaram pelos filtros obrigatórios. Não deve afirmar que examinou todos os perfis do catálogo. Se nenhum dos cinco satisfaz os critérios, a resposta apropriada é registrar a insuficiência ou pedir decisão humana, em vez de escolher uma opção incompatível.

\section[Glossário do capítulo]{Glossário para consultar durante a leitura}

\begin{resultado}
{\footnotesize
\begin{tabularx}{\textwidth}{lY}
\toprule
Termo & Significado neste capítulo \\
\midrule
Agente & Perfil de trabalho com instruções e capacidades declaradas. \\
Blackboard & Quadro compartilhado em que componentes publicam e consultam tarefas. \\
Cartão & Registro estruturado que anuncia a identidade e as capacidades de um agente. \\
Catálogo & Conjunto de fichas originais que descrevem os perfis de agente. \\
Compilador & Programa que lê as fichas e produz o arquivo de configuração. \\
Configuração & Opções que informam ao aplicativo quais recursos estão declarados. \\
Disponível & Recurso que o processo atual consegue carregar e usar. \\
Frontmatter & Bloco no início da ficha que armazena campos como nome, versão e descrição. \\
Metadados & Informações que descrevem outro registro, arquivo ou recurso. \\
Orquestrador & Componente que coordena tarefas e escolhe ou acompanha uma rota. \\
Especificação (spec) & Documento que registra um comportamento esperado e critérios para avaliá-lo. \\
Execução & Trabalho que o programa realmente realizou em uma sessão. \\
Histórico de evolução & Registros de mudanças, com identificadores e contexto. \\
Evento & Ocorrência, como uma tarefa publicada, que pode acionar outro componente. \\
Hook & Programa automático acionado por um evento definido. \\
JSON & Formato de texto estruturado usado para representar dados e opções. \\
MCP & Protocolo Model Context Protocol, pelo qual um aplicativo solicita ferramentas oferecidas por um servidor. \\
MetaBus & Canal local que distribui mensagens e registros entre componentes. \\
MCI & Memória metacognitiva e interconexão: registros sobre trabalho anterior e comunicação. \\
Plugin & Complemento de código que amplia as funções reconhecidas pelo aplicativo. \\
Provider & Componente que conecta o aplicativo a um serviço de modelo. \\
Runtime & Programa enquanto está em execução, com os recursos que conseguiu carregar. \\
Skill & Instruções reutilizáveis que explicam como conduzir um tipo de tarefa. \\
Smoke test & Teste rápido que procura erros comuns em um caminho pequeno de execução. \\
Stub & Substituto simples que imita parte de outro sistema para permitir um teste local. \\
Teste & Execução de uma verificação que compara um resultado com uma condição esperada. \\
Token Economy & Contabilidade interna de unidades associadas às tarefas do sistema. \\
\bottomrule
\end{tabularx}}
\end{resultado}
